\chapter{Future Work}
\label{c:futurework}

The most pressing extension of the present work is the generalisation to
multi-source detection. All architectures evaluated in this thesis assume that
a single patient-zero initiated the outbreak, an assumption that is often
violated during the early stages of an emerging epidemic when multiple
individuals may independently introduce the pathogen into the contact network.
A principled multi-source formulation requires either a multi-label output head
that assigns independent probabilities to each node in a candidate source set
$\mathcal{S}^*$, or a sequential detection scheme in which sources are
identified and removed iteratively, analogous to peeling algorithms in
combinatorial inference. Either approach demands a revised training procedure
and evaluation protocol, since standard top-$k$ accuracy metrics conflate
accuracy with set-size sensitivity when $|\mathcal{S}^*| > 1$. Alongside the
multi-source extension, the inductive generalisation capacity of all evaluated
models deserves attention: currently, models are trained and evaluated on the
same network topology, precluding deployment on previously unseen contact
networks without retraining. Integrating per-node memory modules in the style
of Temporal Graph Networks (TGN) \cite{Rossi2020}---which maintain a
continuously updated state vector for each node and update it via a gated
recurrent unit as new contacts arrive---would enable a model to adapt to
arbitrary network topologies at inference time, a property that is essential
for real-time outbreak response applications where the contact network may
change between training and deployment.

On the architectural side, the De Bruijn GNN can be extended by incorporating
the Hypergeometric Pattern Analysis (HYPA) framework of
\cite{vonPichowski2024}. HYPA uses a hypergeometric ensemble of randomised
temporal networks as a null model and retains only those higher-order temporal
patterns---represented as De Bruijn graph nodes---whose observed frequency is
statistically anomalous relative to the null model. Using HYPA-filtered De
Bruijn graphs as input to the GNN introduces a principled, data-driven inductive
bias that suppresses spurious higher-order patterns arising from random
fluctuations in the contact sequence, potentially improving both accuracy and
generalisation across network families. Beyond epidemic applications, the
general framework developed in this thesis---temporal network representation,
Monte Carlo simulation-based training, and causal GNN inference---transfers
naturally to other spreading processes on temporal graphs, including rumour
and misinformation propagation in social networks, malware transmission in
computer infrastructure, and cascading failures in power or transport systems.
Adapting the simulation engine and loss functions to these domains would
broaden the empirical scope of the methodology considerably. A further
promising direction is the integration of active querying strategies, in which
the model adaptively selects a small number of additional nodes whose infection
status is queried to resolve ambiguity among the top-ranked source candidates
\cite{Sterchi2025}. Such a strategy converts the purely passive inference
problem studied here into a sequential decision problem and offers a principled
way to improve accuracy under tight observation budgets---a practically relevant
setting for epidemiological investigation, where testing resources are limited.
